% Lliçó: el teorema de Bolzano (valencià).
%
% Teorema, demostració, corol·lari i exemple. La demostració va en
% `notesonly`: en les diapositives s'enuncia i en els apunts es demostra.

\begin{frame}
\didactatitle{El teorema de Bolzano}

\begin{theorem}[Bolzano]
Si $f$ és contínua en $[a, b]$ i $f(a)$ i $f(b)$ tenen signes oposats,
existeix $c \in (a, b)$ tal que $f(c) = 0$.
\end{theorem}

\bymedium{%
  Una funció contínua que comença per davall de l'eix i acaba per damunt
  l'ha de travessar.
}{%
  Una funció contínua que comença per davall de l'eix i acaba per damunt
  l'ha de travessar. El que el teorema afig a la intuïció és que això no
  depén del dibuix, sinó que la recta real no té forats: l'axioma del
  suprem.
}

\begin{notesonly}
\begin{proof}
Suposem $f(a) < 0 < f(b)$; l'altre cas es redueix a aquest canviant $f$ per
$-f$. El conjunt $A = \{x \in [a, b] : f(x) < 0\}$ conté $a$ i està fitat
per $b$, així que té suprem $c \in [a, b]$. Vegem que $f(c) = 0$.

Si fóra $f(c) \neq 0$, prenent $\varepsilon = |f(c)|$ en la definició de
continuïtat hi hauria un $\delta > 0$ tal que $f(x)$ té el mateix signe que
$f(c)$ per a tot $x \in [a, b]$ amb $|x - c| < \delta$. Si $f(c) < 0$,
aleshores $c < b$ i hi hauria punts de $A$ més grans que $c$. Si $f(c) > 0$,
aleshores $c > a$ i cap punt de $(c - \delta, c]$ no estaria en $A$, així que
$c - \delta$ seria una fita superior de $A$ menor que $c$. Les dues coses
contradiuen que $c$ siga el suprem.

Per tant $f(c) = 0$, i $c \neq a, b$ perquè ni $f(a)$ ni $f(b)$ s'anul·len.
\end{proof}
\end{notesonly}

\begin{teaching}[Per a discutir]
Què falla si $f$ no és contínua? (El salt $H(x) - \tfrac12$ en $[-1, 1]$.)
I si només existiren els racionals? ($x^2 - 2$ en $[0, 2]$ canvia de signe i
no té cap arrel racional.)
\end{teaching}
\end{frame}

\begin{frame}
\didactatitle{Valors intermedis}

\begin{corollary}[Valors intermedis]
Si $f$ és contínua en $[a, b]$, pren tots els valors compresos entre
$f(a)$ i $f(b)$.
\end{corollary}

\onlynotes{%
  Si $k$ està estrictament entre $f(a)$ i $f(b)$, n'hi ha prou d'aplicar el
  teorema de Bolzano a $g(x) = f(x) - k$, que és contínua i canvia de signe en
  $[a, b]$.
}

\dpause

\begin{example}
L'equació $x^3 + x - 1 = 0$ té una solució en $(0, 1)$: el polinomi és
continu, val $-1$ en $0$ i $1$ en $1$.
\onlynotes{%
  Repetint l'argument en la meitat de l'interval on canvia el signe, i
  després en la meitat d'eixa meitat, s'obté el \keyterm{mètode de
  bisecció}: cada pas divideix per dos la longitud de l'interval en què està
  l'arrel.
}
\end{example}
\end{frame}
